\section{The boundary between the main proof and stochastic analysis}
\label{sec:graph}
The main argument receives its analytic results through the seven propositions in this section.
Each clause identifier corresponds to one theorem replaced by an assumption in the Lean isolation check.
There are twenty clauses, together with two data definitions for the integral graph and the market.
These are the complete inputs to the main proof when the analytic implementation is removed.
Accepting them suffices to read the entire main argument, including Section~\ref{sec:selection} onwards,
without consulting internal appendix lemmas.
Ordinary measure theory, conditional expectation, Doob inequalities, optional sampling and Hilbert-space convex compactness are shared foundations.
NFLVR, maximality, the selection of convex averages, boundedness in probability of whole families,
and the M/FV Cauchy properties are not inputs here: they are used or proved in the main argument.

\subsection{Two data definitions and notation}
Fix the original bounded price $S$ and $(\Omega,\F,(\F_t),\mu)$ satisfying the usual conditions.
Whenever $Q\sim\mu$ is used below, $Q$ is a probability measure and the same filtration satisfies the usual conditions under $Q$.
Unless stated otherwise, process identities and all-time bounds are understood up to indistinguishability or on one common full-measure set.
An elementary predictable test $J$, with $|J|\leq1$, consists of finitely many stopping intervals and coefficients measurable at their left endpoints.
Define $J\cdot Y$ by the corresponding finite sum of increments and write
\[
 e_T^J(Y,Z)=1\wedge\sup_{t\leq T}|(J\cdot(Y-Z))_t|,
 \qquad d_T^Q(Y,Z)=\sup_J E_Q e_T^J(Y,Z).
\]
Emery convergence means $d_T^Q(Y_n,Y)\to0$ for every finite $T$.
The Cauchy quantifiers are $\forall T,\varepsilon>0\ \exists N\ \forall n,k\geq N\ \forall J$.

\begin{definition}[The original-price integral graph and terminal market]\label{def:boundary-data}
\leavevmode\par
\begin{enumerate}
\item \analyticcontract{G0}{FTAPTheorem42.BoundedSourceIntegralMarket.GeneralIntegralGraph}
Put $H^{(n)}=H1_{\{|H|\leq n+1\}}$.
A local special integral is obtained by applying the \emph{same} predictable coefficient to the
energy integral of the M component and the Stieltjes integral of the FV component on stopped coordinates, then gluing.
Write $(H,X)\in\mathcal G_S$, or $X=H\cdot S$, when the local special integrals of $H^{(n)}$
converge in Emery topology to $X$, which has a zero-start strongly adapted c\`adl\`ag representative.
Neither specialness nor existence of a terminal value is required of $X$ itself.
\item \analyticcontract{M0}{FTAPTheorem42.BoundedSourceIntegralMarket.generalMarket}
The gain-process model $\mathcal M_S$ is obtained from the Emery completion of elementary integrals of the original price.
Its strategies carry a realized gain and its own terminal value; gain and terminal value respect linear operations.
Clauses \contractref{M1}--\contractref{M3} below identify this model with the market defined by $\mathcal G_S$.
\end{enumerate}
Set $K_a=\{f:a>0,\ \exists(H,X)\in\mathcal G_S,\ X\geq-a,\ X_t\to f\ \as\}$ and
$K_0=\bigcup_{a>0}K_a$. These set definitions remain concrete in the isolation check.
An ``original-market gain'' means a zero-start strongly adapted c\`adl\`ag process $Y$ with a witness $(H,Y)\in\mathcal G_S$.
A ``special gain'' additionally carries a specified zero-start decomposition $Y=M+A$ under $Q$,
where $M$ is a c\`adl\`ag local martingale and $A$ is predictable, c\`adl\`ag and locally of finite variation.
A ``finite gain'' is a special gain constant after some finite time, whose specified $A$ is globally BV.
These data types also remain in the isolation check.
\end{definition}

\subsection{Seven analytic inputs}
\begin{proposition}[Regular limits, common envelopes and equivalent measures]\label{input:paths}
\leavevmode\par
\begin{enumerate}
\item \analyticcontract{P1}{FTAPTheorem42.AnalyticInterface.regular_uniform_limit}
If strongly adapted c\`adl\`ag processes $Y_n$ are Cauchy in probability for the all-time supremum,
there are a strictly increasing $f$ and a strongly adapted c\`adl\`ag $X$ such that $Y_{f(n)}\to X$ uniformly over all time almost surely.
\item \analyticcontract{P2}{FTAPTheorem42.AnalyticInterface.regular_common_envelope}
If each process in such a Cauchy sequence has its own finite terminal value, there are a strictly increasing $f$
and a finite nonnegative measurable $q$ such that $|Y_{f(n)}(t)|\leq q$ for every $n,t$, almost surely.
\item \analyticcontract{P3}{FTAPTheorem42.AnalyticInterface.equivalent_envelope_measure}
For any finite nonnegative measurable $q$, there is an equivalent probability $Q\sim\mu$ with $q\in L^2(Q)$;
the same filtration satisfies the usual conditions under $Q$.
\end{enumerate}
\end{proposition}

\begin{proposition}[Special decomposition of a specified gain]\label{input:decomposition}
\analyticcontract{D1}{FTAPTheorem42.BoundedSourceIntegralMarket.originalGain_specialDecomposition}
If an original-market gain $Y$ satisfies $|Y_t|\leq|q|$ for $q\in L^2(Q)$,
the same $Y$ has a zero-start special decomposition $Y=M+A$.
The decomposition is constructed under $Q$; no transfer of specialness to the original measure is asserted.
\end{proposition}

\begin{proposition}[Market identification and finite operations]\label{input:market}
\leavevmode\par
\begin{enumerate}
\item \analyticcontract{M1}{FTAPTheorem42.BoundedSourceIntegralMarket.generalAdmissibleClaims_eq_generalMarket}
For every real $a$, $K_a$ equals the $a$-admissible terminal set of $\mathcal M_S$.
\item \analyticcontract{M2}{FTAPTheorem42.BoundedSourceIntegralMarket.generalTerminalClaims_eq_generalMarket_K0}
$K_0$ equals the admissible terminal set of $\mathcal M_S$.
\item \analyticcontract{M3}{FTAPTheorem42.BoundedSourceIntegralMarket.generalAdmissibleClaims_one_eq_generalMarket_K1}
$K_1$ equals the unit-admissible terminal set of $\mathcal M_S$.
\item \analyticcontract{M4}{FTAPTheorem42.BoundedSourceIntegralMarket.generalMarket_terminal_properties}
Admissible terminal claims of this model are a.e. strongly measurable, and an admissibility bound $-a$ also holds at the terminal value.
\item \analyticcontract{M5}{FTAPTheorem42.BoundedSourceIntegralMarket.generalMarket_operations}
Model gains are zero-start, strongly adapted and c\`adl\`ag, have their own terminal values, and are realized in $\mathcal G_S$.
An a.e. equal terminal representative can be substituted without changing the gain.
For gains $Y,Z$, a bounded stopping time $\tau\leq T$, and $E\in\F_\tau$, there is a switched strategy
with gain $Y_{t\wedge\tau}+Z_t-Z_{t\wedge\tau}$ on $E$, and $Y_t$ on $E^c$.
Its terminal value is $Z_\infty+Y_\tau-Z_\tau$ on $E$, and $Y_\infty$ on $E^c$.
\item \analyticcontract{M6}{FTAPTheorem42.BoundedSourceIntegralMarket.originalSpecialGain_finiteStop}
For a special gain $Y=M+A$ and finite $T>0$, $Y^T=M^T+A^T$ is realized as a finite gain.
\item \analyticcontract{M7}{FTAPTheorem42.BoundedSourceIntegralMarket.originalGain_convexCombination}
Finite convex combinations of special gains are realized using the \emph{same weights} for gains, M components and FV components.
In particular this realization is available for specified forward convex weights.
\end{enumerate}
\end{proposition}

\begin{proposition}[Finite-scale quantitative estimates]\label{input:finite}
None of the following clauses assumes NFLVR.
\leavevmode\par
\begin{enumerate}
\item \analyticcontract{F1}{FTAPTheorem42.BoundedSourceIntegralMarket.originalGain_finiteRisk}
Suppose a special gain $Y=M+A\geq-1$ satisfies $|Y|\leq\xi$ and $\|\xi\|_2\leq B$.
Let $0<\alpha\leq1/4$, let $n$ be a positive integer, and let $T>0$, with
\[
 k=\lfloor\alpha n/4\rfloor>0,\quad 6B\leq n^2,\quad (B/n)^2\leq\alpha,
 \quad m=\alpha^2-(4B/(\alpha n^2))^2\geq0.
\]
If $Q(\sup_{t\leq T}|M_t|>n^3)>8\alpha$, there is an a.e. measurable $f\in K_0-L^0_+(Q)$ such that
\[
 f\geq-n^{-1/4}-\frac2{nk},\qquad
 Q\left(f\geq\frac{\alpha m}4-n^{-1/4}\right)
 \geq\frac m2-\frac{16\sqrt n}{k}.
\]
\item \analyticcontract{F2}{FTAPTheorem42.BoundedSourceIntegralMarket.originalGain_finiteTail_normalization}
Let finite gains $Y_i=M_i+A_i\geq-1$ have a common envelope $|q|\in L^2(Q)$ and suppose $\{M_i^*\}$ is bounded in probability.
Put $\tau_i(c)=\inf\{t:|M_i(t)|>c\}$ and $L^{w,c}=\sum_iw_i(M_i-M_i^{\tau_i(c)})$.
Let $B(w,c,\varepsilon)$ be the event that the passage time of $|L^{w,c}|$ across level $\varepsilon$ is finite.
For every $\varepsilon>0$ there is $N\geq0$ such that for every $0<\delta\leq\varepsilon/4$ there is $c_0\geq0$ with the following property.
For every finite convex weight $w$ and $c\geq c_0$, if $Q(B(w,c,\varepsilon))>\varepsilon$,
there are an original-market special gain $V=M^V+A^V$ and an envelope $g$ with
\[
 V\geq-1,\quad |V|\leq g,\quad \|g\|_2\leq1,\quad
 Q((M^V)^*>\varepsilon/[2(1+N)\delta])>\varepsilon/2.
\]
\item \analyticcontract{F3}{FTAPTheorem42.BoundedSourceIntegralMarket.originalGain_finiteTail_test_estimate}
Assume the finite-gain, common-envelope and M-maximal boundedness conditions of the preceding clause, but not the lower bound $-1$.
For every $\eta>0$ there is $c_0\geq0$ such that for every $w$, $c\geq c_0$, $T>0$ and $J$,
\[
 E_Q e_T^J(L^{w,c},0)\leq4\eta+Q(B(w,c,\eta)).
\]
\item \analyticcontract{F4}{FTAPTheorem42.BoundedSourceIntegralMarket.originalGain_prefix_energy}
For a special gain $Y=M+A$, $q\in L^2(Q)$, $|Y|\leq q$, $c\geq0$ and $T>0$,
$P=M^{\tau(c)\wedge T}$ is a true martingale, $P_T\in L^2(Q)$, and
$\|P_T\|_2\leq c+6\|q\|_2$.
\item \analyticcontract{F5}{FTAPTheorem42.BoundedSourceIntegralMarket.originalGain_prefix_regular}
The same stopped component $P=M^{\tau(c)\wedge T}$ is right-continuous and satisfies $P_0=0$ almost surely
for every real $c$ and $T\geq0$, without an envelope assumption.
\end{enumerate}
\end{proposition}

\begin{proposition}[Martingale approximation uniformly over tests]\label{input:approximation}
\leavevmode\par
\begin{enumerate}
\item \analyticcontract{A1}{FTAPTheorem42.AnalyticInterface.emeryCauchy_of_uniform_approximations}
Let $X_n$ be strongly progressive. Suppose for every $T,\varepsilon>0$ there is a strongly progressive sequence $P_n$
with $\sup_n d_T^Q(X_n,P_n)\leq\varepsilon$, Cauchy uniformly over tests at this $T$.
Then $X_n$ is Emery Cauchy.
\item \analyticcontract{A2}{FTAPTheorem42.AnalyticInterface.martingaleTestsCauchy_of_terminalL2}
Let $P_n$ be right-continuous zero-start true martingales and let $U\geq0$ be finite.
If $(P_n(U))$ is Cauchy in $L^2(Q)$, then for every $\varepsilon>0$ there is $N$ such that
$E_Qe_T^J(P_n,P_k)\leq\varepsilon$ for every $n,k\geq N$, every $T\leq U$, and every $J$.
\end{enumerate}
\end{proposition}

\begin{proposition}[Finite Hahn trades for two specified gains]\label{input:hahn}
\analyticcontract{H1}{FTAPTheorem42.BoundedSourceIntegralMarket.originalGain_finiteHahn}
Let special gains $Y=M_Y+A_Y$ and $Z=M_Z+A_Z$ share an envelope $q\in L^2(Q)$ and both be bounded below by $-1$.
Fix $T>0$ and $\delta\geq0$, and assume $E_Qe_T^J(M_Y,M_Z)\leq b$ and the reverse-order estimate for every $J$.
The following data can be chosen simultaneously for both orientations. For orientation $(Y,Z)$ put $L=M_Y-M_Z$ and $A=A_Y-A_Z$.
There are strongly adapted $N,C$ and $V=N+C$, with $N$ right-continuous, $C_0=0$, and
\[
 0\leq C_c-C_a,\quad A_{c\wedge T}-A_{a\wedge T}\leq C_c-C_a\quad(a\leq c),
 \qquad E_Q[1\wedge\sup_{t\leq T}|N_t|]\leq b.
\]
Put $\tau=\inf\{t:N_t-\max(L_t,0)<-\delta\}$, $\sigma=\tau\wedge T$, and $E=\{T<\tau\}$.
If $0<\delta\leq1$, then $Q(\tau\leq T)\leq8b/\delta$; for every $U\geq T$,
\[
 Z^{\rm paste}_t=(Z+V)_{t\wedge\sigma}
       +1_E\bigl(Z_{t\wedge U}-Z_{t\wedge T}\bigr),\qquad
 Z^{\rm paste}_U\in K_{1+\delta}.
\]
Writing $C'$ for the FV component in the reverse orientation, for every real $\alpha$,
\[
 Q(\Var_{[0,T]}A>\alpha)>\alpha
 \ \Longrightarrow\ Q(C_T>\alpha/2)>\alpha/2
 \ \text{or}\ Q(C'_T>\alpha/2)>\alpha/2.
\]
This is a finite construction for two gains; it assumes neither a Cauchy sequence, NFLVR nor maximality.
\end{proposition}

\begin{proposition}[Realizing the intrinsic terminal value from component estimates]\label{input:realization}
\analyticcontract{C1}{FTAPTheorem42.BoundedSourceIntegralMarket.truncatedTerminalClaims_one_of_component_estimates}
Let $(H_n,Y_n)\in\mathcal G_S$, with $Y_n$ strongly adapted and right-continuous.
Suppose $Y_n$ converges almost surely uniformly over all time to a zero-start strongly adapted c\`adl\`ag $X\geq-1$, and $X_t\to h$ almost surely.
Under $Q\sim\mu$, let the same gains satisfy $Y_n=M_n+A_n$, where each $A_n$ is right-continuous and every difference $A_m-A_n$ is locally BV on every path.
If for every finite $T$ and $\varepsilon>0$ there is $N$ such that for every $m,n\geq N$ and every $J$,
\[
 Q(e_T^J(M_m,M_n)>\varepsilon)\leq\varepsilon,
 \qquad Q(\Var_{[0,T]}(A_m-A_n)>\varepsilon)\leq\varepsilon,
\]
then $h\in K_1$. No additional martingale assumption on $M_n$ is imposed here.
\end{proposition}

\subsection{Proof locations and the isolation check}
Appendix~\ref{sec:input-proofs} assembles the proofs of the seven inputs.
\begin{center}\small
\begin{tabular}{ll}
\toprule Input & Proof location\\\midrule
Proposition~\ref{input:paths} & Section~\ref{sec:path-inputs}\\
Propositions~\ref{input:decomposition}, \ref{input:market} & Sections~\ref{sec:integral-background}--\ref{sec:auxiliary-transforms}\\
Proposition~\ref{input:finite} & Sections~\ref{sec:ds-estimates}, \ref{sec:approximation-inputs}\\
Proposition~\ref{input:approximation} & Section~\ref{sec:approximation-inputs}\\
Proposition~\ref{input:hahn} & Section~\ref{sec:finite-hahn-proof}\\
Proposition~\ref{input:realization} & Section~\ref{sec:component-realization}\\\bottomrule
\end{tabular}\end{center}
Appendices~\ref{sec:construction}--\ref{sec:emery-realization} divide the internal construction into three stages.
Their internal theorems need not be held in mind simultaneously to follow the main proof.
The Lean correspondence table gives the fully qualified name for each identifier.
The isolation check compares the Japanese and English mappings with the actual set of 22 declarations replaced by assumptions,
and rebuilds the unchanged main proof without \path{Stochastic/}.
It also verifies that the resulting main theorem uses all 22 specified inputs and no other analytic axioms.
This verifies the formal dependency boundary. Agreement in mathematical meaning between the prose and the Lean types
still requires comparing the hypotheses and conclusions stated here with those types.
